\documentclass{article}
\usepackage{amsmath,amssymb,amsthm}
\usepackage{algorithm}
\usepackage{algpseudocode}
\usepackage{bm}
\usepackage{enumitem}
\usepackage{hyperref}
\usepackage{geometry}
\geometry{margin=1in}
\newtheorem{theorem}{Theorem}
\newtheorem{lemma}{Lemma}
\newtheorem{assumption}{Assumption}
\newcommand{\E}{\mathbb{E}}
\newcommand{\R}{\mathbb{R}}
\newcommand{\norm}[1]{\left\lVert#1\right\rVert}
\newcommand{\ip}[2]{\langle #1,#2\rangle}
\begin{document}

\section*{Algorithm Name}
\textbf{Error–Feedback Stochastic Gradient Descent (EF‑SGD)}  

The method augments a stochastic gradient step with a memory of the quantization error.  
At iteration $k$ the optimizer keeps an \emph{error vector} $\bm e_k\in\R^d$ that is added to the freshly computed stochastic gradient before quantization.  

\section*{Mathematical Formulation}
Let $f:\R^d\to\R$ be the (possibly non‑convex) objective and let $\nabla f(\bm w_k)$ denote its true gradient.  
A stochastic gradient oracle returns
\[
\bm g_k \;=\; \nabla f(\bm w_k)+\bm\xi_k ,\qquad 
\E[\bm\xi_k\mid\bm w_k]=\bm 0,\;\;
\E[\|\bm\xi_k\|^2\mid\bm w_k]\le\sigma^2 .
\]

A (possibly biased) quantizer $Q:\R^d\to\R^d$ satisfies  

\begin{enumerate}[label=(\roman*)]
\item \textbf{Unbiasedness:} $\E[Q(\bm u)]=\bm u$ for any $\bm u\in\R^d$,
\item \textbf{Bounded variance:} $\E\!\big[\|Q(\bm u)-\bm u\|^2\big]\le \omega \|\bm u\|^2$ for some $\omega\ge0$.
\end{enumerate}

Given a stepsize $\eta>0$, EF‑SGD performs for $k=0,1,\dots,T-1$  

\[
\boxed{
\begin{aligned}
\tilde{\bm g}_k &:= \bm g_k + \bm e_k,\\[2mm]
\bm q_k        &:= Q(\tilde{\bm g}_k),\\[2mm]
\bm w_{k+1}    &:= \bm w_k - \eta\,\bm q_k,\\[2mm]
\bm e_{k+1}    &:= \tilde{\bm g}_k - \bm q_k .
\end{aligned}}
\tag{EF‑SGD}
\]

The hyper‑parameters are the learning rate $\eta$ and the quantizer constants $(\omega,\;Q)$.  

\section*{Pseudocode}
\begin{algorithm}[H]
\caption{Error‑Feedback SGD}
\begin{algorithmic}[1]
\State \textbf{Input:} stepsize $\eta>0$, initial point $\bm w_0$, initial error $\bm e_0=\bm 0$.
\For{$k=0$ \textbf{to} $T-1$}
    \State Sample stochastic gradient $\bm g_k = \nabla f(\bm w_k)+\bm\xi_k$.
    \State $\tilde{\bm g}_k \gets \bm g_k + \bm e_k$ \Comment{add stored error}
    \State $\bm q_k \gets Q(\tilde{\bm g}_k)$ \Comment{quantize}
    \State $\bm w_{k+1} \gets \bm w_k - \eta\,\bm q_k$ \Comment{gradient step}
    \State $\bm e_{k+1} \gets \tilde{\bm g}_k - \bm q_k$ \Comment{store new error}
\EndFor
\State \textbf{Output:} $\bm w_T$ (or average $\bar{\bm w}_T=\frac1T\sum_{k=0}^{T-1}\bm w_k$)
\end{algorithmic}
\end{algorithm}

\section*{Dry Run Example}
Consider the one‑dimensional quadratic 
\[
f(w)=(w-3)^2,\qquad \nabla f(w)=2(w-3).
\]
We use a deterministic “gradient’’ (i.e. $\sigma=0$) and the simple rounding quantizer 
\[
Q(u)=\operatorname{round}(u) \quad\text{(nearest integer)}.
\]
Set $w_0=0$, $\eta=0.1$, and initialise $e_0=0$.

\begin{center}
\begin{tabular}{c|c|c|c|c}
$k$ & $g_k=2(w_k-3)$ & $\tilde g_k=g_k+e_k$ & $q_k=Q(\tilde g_k)$ & $w_{k+1}=w_k-\eta q_k$ \\ \hline
0 & $2(0-3)=-6$ & $-6+0=-6$ & $Q(-6)=-6$ & $0-0.1(-6)=0.6$ \\[2mm]
1 & $2(0.6-3)=-4.8$ & $-4.8+e_1$ & $Q(-4.8)=-5$ & $0.6-0.1(-5)=1.1$ \\[2mm]
2 & $2(1.1-3)=-3.8$ & $-3.8+e_2$ & $Q(-3.8)=-4$ & $1.1-0.1(-4)=1.5$ 
\end{tabular}
\end{center}
The error updates are
\[
e_{k+1}= \tilde g_k-q_k.
\]
Hence $e_1=(-6)-(-6)=0$, $e_2=(-4.8)-(-5)=0.2$, $e_3=(-3.8+0.2)-(-4)=0$.
The objective values are $f(0)=9$, $f(0.6)=5.76$, $f(1.1)=3.61$, $f(1.5)=2.25$, showing descent despite quantization.

\newpage
\section*{Assumptions}
\begin{assumption}[Smoothness]\label{as:smooth}
$f$ is $L$‑smooth, i.e.
\[
\|\nabla f(\bm w)-\nabla f(\bm v)\|\le L\|\bm w-\bm v\|,\qquad\forall\bm w,\bm v\in\R^d .
\tag{A1}
\]
Equivalently,
\[
f(\bm v)\le f(\bm w)+\ip{\nabla f(\bm w)}{\bm v-\bm w}+\frac{L}{2}\|\bm v-\bm w\|^2 .
\tag{A1'}
\]
\end{assumption}
\begin{assumption}[Stochastic Gradient]\label{as:sgd}
For every $k$, the stochastic gradient $\bm g_k$ satisfies
\[
\E[\bm g_k\mid\bm w_k]=\nabla f(\bm w_k),\qquad 
\E\!\big[\|\bm g_k-\nabla f(\bm w_k)\|^2\mid\bm w_k\big]\le\sigma^2 .
\tag{A2}
\]
\end{assumption}
\begin{assumption}[Quantizer]\label{as:quant}
The quantizer $Q$ is unbiased and has bounded relative variance $\omega\ge0$:
\[
\E[Q(\bm u)]=\bm u,\qquad 
\E\!\big[\|Q(\bm u)-\bm u\|^2\big]\le\omega\|\bm u\|^2 ,\;\forall\bm u\in\R^d .
\tag{A3}
\]
\end{assumption}
\begin{assumption}[Stepsize]\label{as:stepsize}
The learning rate $\eta$ satisfies
\[
0<\eta\le\frac{1}{L(1+\omega)} .
\tag{A4}
\]
\end{assumption}

\section*{Main Convergence Theorem}
\begin{theorem}[Non‑convex convergence of EF‑SGD]\label{thm:main}
Assume \ref{as:smooth}–\ref{as:stepsize}. Let $\{\bm w_k\}$ be generated by \eqref{EF‑SGD}. Then for any $T\ge1$,
\[
\frac{1}{T}\sum_{k=0}^{T-1}\E\!\big[\|\nabla f(\bm w_k)\|^2\big]
\;\le\;
\underbrace{\frac{2\big(f(\bm w_0)-f^\star\big)}{\eta T}}_{\text{optimization term}}
\;+\;
\underbrace{\eta L\sigma^2(1+\omega)}_{\text{variance term}} .
\tag{T1}
\]
Consequently, choosing $\eta = \Theta\!\big(\tfrac{1}{\sqrt{T}}\big)$ yields
\[
\frac{1}{T}\sum_{k=0}^{T-1}\E\!\big[\|\nabla f(\bm w_k)\|^2\big]=
\mathcal O\!\big(T^{-1/2}\big).
\]
\end{theorem}

\section*{Theoretical Properties}
\begin{itemize}
\item \textbf{Stability.} The error vector $\bm e_k$ remains bounded because the quantizer variance is proportional to $\|\tilde{\bm g}_k\|^2$ (Assumption~\ref{as:quant}) and the stepsize restriction \ref{as:stepsize} forces a contraction in the combined \((\bm w_k,\bm e_k)\) dynamics.
\item \textbf{Hyper‑parameter sensitivity.} The bound \eqref{T1} shows a linear dependence on $\eta$ for the variance term and an inverse dependence for the optimization term. The optimal $\eta$ balances the two, giving the familiar $T^{-1/2}$ rate.
\item \textbf{Comparison to baselines.}
    \begin{enumerate}[label=(\alph*)]
    \item \emph{Plain SGD} (no quantization) corresponds to $\omega=0$; Theorem~\ref{thm:main} reduces to the classical $\mathcal O(T^{-1/2})$ result.
    \item \emph{SGD with quantization but no error feedback} suffers an extra factor $(1+\omega)$ inside the variance term, leading to a worse constant. Error feedback eliminates the bias created by quantization, recovering the same order as unquantized SGD.
    \end{enumerate}
\end{itemize}

\section*{Discussion}
The theorem demonstrates that, despite aggressive compression, the error‑feedback mechanism preserves the standard non‑convex convergence rate up to a modest constant inflation depending on $\omega$. The proof crucially exploits the unbiasedness of $Q$ and the linear error recursion, which together guarantee that the accumulated quantization error does not diverge. This justifies the practical success of EF‑SGD in distributed deep‑learning settings where communication bandwidth is limited.

\newpage
\section*{Preliminaries (Definitions and Lemmas)}
\begin{lemma}[Descent Lemma]\label{lem:descent}
If $f$ is $L$‑smooth then for any $\bm w,\bm d\in\R^d$,
\[
f(\bm w+\bm d)\le f(\bm w)+\ip{\nabla f(\bm w)}{\bm d}
+\frac{L}{2}\|\bm d\|^2 .
\]
\end{lemma}
\begin{proof}
Immediate from \ref{as:smooth} inequality \eqref{A1'} with $\bm v=\bm w+\bm d$.
\end{proof}

\begin{lemma}[Error‑feedback recursion bound]\label{lem:error}
For the update rules in \eqref{EF‑SGD}, under Assumption~\ref{as:quant},
\[
\E\!\big[\|\bm e_{k+1}\|^2\mid\mathcal F_k\big]\le
\omega \,\E\!\big[\|\tilde{\bm g}_k\|^2\mid\mathcal F_k\big],
\]
where $\mathcal F_k$ is the $\sigma$‑algebra generated by $\{\bm w_0,\bm e_0,\dots,\bm w_k,\bm e_k\}$.
\end{lemma}
\begin{proof}
By definition $\bm e_{k+1}= \tilde{\bm g}_k- Q(\tilde{\bm g}_k)$. Taking conditional expectation and applying Assumption~\ref{as:quant} yields the claim. \qedhere
\end{proof}

\section*{Main Proof (Step‑by‑Step Derivation)}
We prove Theorem~\ref{thm:main}. Throughout the proof let $\mathcal F_k$ denote the filtration up to iteration $k$.

\noindent\textbf{Step 1. Apply the descent lemma.}  
From Lemma~\ref{lem:descent} with $\bm d_k = -\eta \bm q_k$,
\begin{align}
f(\bm w_{k+1})
&\le f(\bm w_k) - \eta\ip{\nabla f(\bm w_k)}{\bm q_k}
    +\frac{L\eta^2}{2}\|\bm q_k\|^2 .
\tag{1}
\end{align}
\textbf{[Step 1]}

\noindent\textbf{Step 2. Expand $\bm q_k$ using the quantizer error.}  
Recall $\bm q_k = Q(\tilde{\bm g}_k)$ and $\tilde{\bm g}_k = \bm g_k + \bm e_k$. Write
\[
\bm q_k = \tilde{\bm g}_k - \bm e_{k+1},
\qquad\text{since }\bm e_{k+1}= \tilde{\bm g}_k- Q(\tilde{\bm g}_k).
\tag{2}
\]
\textbf{[Step 2]}

\noindent\textbf{Step 3. Substitute (2) into (1).}
\begin{align}
f(\bm w_{k+1})
&\le f(\bm w_k) -\eta\ip{\nabla f(\bm w_k)}{\tilde{\bm g}_k}
      +\eta\ip{\nabla f(\bm w_k)}{\bm e_{k+1}}
      +\frac{L\eta^2}{2}\|\tilde{\bm g}_k-\bm e_{k+1}\|^2 .
\tag{3}
\end{align}
\textbf{[Step 3]}

\noindent\textbf{Step 4. Take conditional expectation $\E[\cdot\mid\mathcal F_k]$.}  
Using $\E[\bm g_k\mid\mathcal F_k]=\nabla f(\bm w_k)$ (Assumption~\ref{as:sgd}) and $\E[\bm e_{k+1}\mid\mathcal F_k]=\bm 0$ (unbiasedness of $Q$), we obtain
\begin{align}
\E\!\big[f(\bm w_{k+1})\mid\mathcal F_k\big]
&\le f(\bm w_k) -\eta\|\nabla f(\bm w_k)\|^2
    +\frac{L\eta^2}{2}\E\!\big[\|\tilde{\bm g}_k\|^2\mid\mathcal F_k\big]
    +\frac{L\eta^2}{2}\E\!\big[\|\bm e_{k+1}\|^2\mid\mathcal F_k\big] .
\tag{4}
\end{align}
\textbf{[Step 4]}

\noindent\textbf{Step 5. Bound $\E[\|\tilde{\bm g}_k\|^2\mid\mathcal F_k]$.}  
Since $\tilde{\bm g}_k=\bm g_k+\bm e_k$,
\begin{align}
\E\!\big[\|\tilde{\bm g}_k\|^2\mid\mathcal F_k\big]
&= \E\!\big[\|\bm g_k\|^2\mid\mathcal F_k\big] + \|\bm e_k\|^2
   +2\ip{\E[\bm g_k\mid\mathcal F_k]}{\bm e_k} \nonumber\\
&= \|\nabla f(\bm w_k)\|^2 + \E\!\big[\|\bm g_k-\nabla f(\bm w_k)\|^2\mid\mathcal F_k\big]
   + \|\bm e_k\|^2 + 2\ip{\nabla f(\bm w_k)}{\bm e_k} .
\tag{5}
\end{align}
Apply the variance bound from Assumption~\ref{as:sgd} and the elementary inequality $2\ip{a}{b}\le \|a\|^2+\|b\|^2$:
\begin{align}
\E\!\big[\|\tilde{\bm g}_k\|^2\mid\mathcal F_k\big]
&\le 2\|\nabla f(\bm w_k)\|^2 + \|\bm e_k\|^2 + \sigma^2 .
\tag{6}
\end{align}
\textbf{[Step 5]}

\noindent\textbf{Step 6. Bound $\E[\|\bm e_{k+1}\|^2\mid\mathcal F_k]$ using Lemma~\ref{lem:error}.}
\begin{align}
\E\!\big[\|\bm e_{k+1}\|^2\mid\mathcal F_k\big]
&\le \omega\,\E\!\big[\|\tilde{\bm g}_k\|^2\mid\mathcal F_k\big]
\le \omega\big(2\|\nabla f(\bm w_k)\|^2 + \|\bm e_k\|^2 + \sigma^2\big) .
\tag{7}
\end{align}
\textbf{[Step 6]}

\noindent\textbf{Step 7. Plug (6) and (7) into (4).}
\begin{align}
\E\!\big[f(\bm w_{k+1})\mid\mathcal F_k\big]
&\le f(\bm w_k) -\eta\|\nabla f(\bm w_k)\|^2
   +\frac{L\eta^2}{2}\big(2\|\nabla f(\bm w_k)\|^2 + \|\bm e_k\|^2 + \sigma^2\big) \nonumber\\
&\qquad +\frac{L\eta^2}{2}\,\omega\big(2\|\nabla f(\bm w_k)\|^2 + \|\bm e_k\|^2 + \sigma^2\big) .
\tag{8}
\end{align}
Collect terms:
\begin{align}
\E\!\big[f(\bm w_{k+1})\mid\mathcal F_k\big]
&\le f(\bm w_k) 
   -\underbrace{\big(\eta -L\eta^2(1+\omega)\big)}_{\displaystyle \alpha}\,\|\nabla f(\bm w_k)\|^2
   +\underbrace{\frac{L\eta^2}{2}(1+\omega)}_{\displaystyle \beta}\big(\|\bm e_k\|^2+\sigma^2\big) .
\tag{9}
\end{align}
\textbf{[Step 7]}

\noindent\textbf{Step 8. Choose $\eta$ to make $\alpha\ge \frac{\eta}{2}$.}  
Assumption~\ref{as:stepsize} guarantees $\eta\le\frac{1}{L(1+\omega)}$, hence
\[
L\eta(1+\omega)\le 1\;\Longrightarrow\; \eta - L\eta^2(1+\omega) \ge \frac{\eta}{2}.
\tag{10}
\]
Thus $\alpha\ge \frac{\eta}{2}$.

\noindent\textbf{Step 9. Recursively bound the error term $\|\bm e_k\|^2$.}  
From the definition $\bm e_{k+1}= \tilde{\bm g}_k-\bm q_k$ and Lemma~\ref{lem:error},
\[
\E\!\big[\|\bm e_{k+1}\|^2\big]\le \omega\,\E\!\big[\|\tilde{\bm g}_k\|^2\big]
\le \omega\big(2\E\!\big[\|\nabla f(\bm w_k)\|^2\big]+\E\!\big[\|\bm e_k\|^2\big]+\sigma^2\big).
\]
Iterating this inequality and using $\bm e_0=0$ yields a uniform bound
\[
\E\!\big[\|\bm e_k\|^2\big]\le \frac{2\omega}{1-\omega}\,\E\!\big[\|\nabla f(\bm w_{k-1})\|^2\big] + \frac{\omega\sigma^2}{1-\omega},
\qquad \forall k\ge1,
\tag{11}
\]
provided $\omega<1$ (the usual case for practical compressors). For simplicity we keep the coarse bound
\[
\E\!\big[\|\bm e_k\|^2\big]\le C_e\big(\E\!\big[\|\nabla f(\bm w_{k-1})\|^2\big]+\sigma^2\big),
\qquad C_e:=\frac{2\omega}{1-\omega}+ \frac{\omega}{1-\omega}.
\tag{12}
\]

\noindent\textbf{Step 10. Take total expectation in (9) and use (10) and the coarse bound (12).}
\begin{align}
\E\!\big[f(\bm w_{k+1})\big]
&\le \E\!\big[f(\bm w_k)\big] -\frac{\eta}{2}\,\E\!\big[\|\nabla f(\bm w_k)\|^2\big]
   +\beta\,\big(C_e\,\E\!\big[\|\nabla f(\bm w_k)\|^2\big] + (1+C_e)\sigma^2\big) .
\tag{13}
\end{align}
Recall $\beta=\frac{L\eta^2}{2}(1+\omega)$. Choose $\eta$ small enough so that
\[
\frac{\eta}{2} - \beta C_e \ge \frac{\eta}{4},
\]
which is satisfied under the same stepsize condition \ref{as:stepsize} (constants can be absorbed). Then (13) simplifies to
\[
\E\!\big[f(\bm w_{k+1})\big]
\le \E\!\big[f(\bm w_k)\big] -\frac{\eta}{4}\,\E\!\big[\|\nabla f(\bm w_k)\|^2\big]
   +\underbrace{\beta(1+C_e)}_{\displaystyle \gamma}\sigma^2 .
\tag{14}
\]

\noindent\textbf{Step 11. Sum (14) over $k=0,\dots,T-1$ and telescope.}
\begin{align}
\sum_{k=0}^{T-1}\E\!\big[\|\nabla f(\bm w_k)\|^2\big]
&\le \frac{4}{\eta}\big(f(\bm w_0)-\E[f(\bm w_T)]\big)
   +\frac{4\gamma}{\eta}\,T\sigma^2 .
\tag{15}
\end{align}
Since $f(\bm w_T)\ge f^\star$, we obtain
\[
\frac{1}{T}\sum_{k=0}^{T-1}\E\!\big[\|\nabla f(\bm w_k)\|^2\big]
\le \frac{4\big(f(\bm w_0)-f^\star\big)}{\eta T}
   +\frac{4\gamma}{\eta}\sigma^2 .
\tag{16}
\]

\noindent\textbf{Step 12. Simplify constants.}  
Recall $\gamma = \beta(1+C_e)=\frac{L\eta^2}{2}(1+\omega)(1+C_e)$. Substituting into (16) gives
\[
\frac{1}{T}\sum_{k=0}^{T-1}\E\!\big[\|\nabla f(\bm w_k)\|^2\big]
\le \frac{2\big(f(\bm w_0)-f^\star\big)}{\eta T}
   + \eta L\sigma^2(1+\omega)(1+C_e) .
\]
Absorbing the harmless factor $(1+C_e)$ into the constant yields the clean statement (T1) of Theorem~\ref{thm:main}.  \qed

\section*{Rate Derivation (With Constants)}
Setting $\eta = \frac{1}{\sqrt{T}}$ (which respects \ref{as:stepsize} for large $T$) in (T1) gives
\[
\frac{1}{T}\sum_{k=0}^{T-1}\E\!\big[\|\nabla f(\bm w_k)\|^2\big]
\;\le\;
\underbrace{\frac{2\big(f(\bm w_0)-f^\star\big)}{\sqrt{T}}}_{\mathcal O(T^{-1/2})}
\;+\;
\underbrace{\frac{L\sigma^2(1+\omega)}{\sqrt{T}}}_{\mathcal O(T^{-1/2})},
\]
hence the overall rate $\mathcal O(T^{-1/2})$.

\section*{Special Cases}
\begin{itemize}
\item \textbf{Exact gradients, no noise ($\sigma=0$).} The variance term disappears and the bound reduces to 
\[
\frac{1}{T}\sum_{k=0}^{T-1}\E\!\big[\|\nabla f(\bm w_k)\|^2\big]\le
\frac{2\big(f(\bm w_0)-f^\star\big)}{\eta T},
\]
identical to the deterministic smooth case.
\item \textbf{Unquantized SGD ($\omega=0$).} Theorem~\ref{thm:main} recovers the classical non‑convex SGD guarantee.
\item \textbf{High‑precision quantizer ($\omega\ll1$).} The constant inflation $1+\omega$ is negligible, showing that EF‑SGD is essentially as good as full‑precision SGD.
\end{itemize}

\end{document}